% Research contribution of /root/h9_height, JAC1964-20261121.
% This is a section fragment, not a standalone document.
% Sources are identified in H9_sources.md. All claims below have an explicit
% logical status. No formal proof certificate is asserted.

\section{The pencil formula: a counterexample and its local repair}

The text of cote 44, scan 4, assumes that the restrictions of the two
cycles are algebraically equivalent to zero. H9 replaces this hypothesis by
homological triviality. These are different hypotheses. Moreover, geometric
integrality of the fibres does not remove higher-codimension vertical
corrections. The following counterexample satisfies even the stronger
condition that the restrictions are rationally equivalent to zero.

\subsection*{H9-C1. A quadric counterexample, with both local contributions}

\textbf{Claim.} For every integer $m\geq2$, over every algebraically closed
field $k$, there is a
Lefschetz pencil on a smooth $2m$-dimensional quadric $X$, all of whose fibres
are geometrically integral, and a codimension-$m$ cycle $Z$ such that
\[
 Z|_{Y}=0\quad\text{in }\operatorname{CH}^{m}(Y),\qquad
 (-1)^m\deg(Z^2)=2.
\]
Here $Y$ is the generic fibre over $K=k(t)$. Consequently every
bilinear pairing on its Chow group gives height zero to the restriction,
whereas the uncorrected expression in H9(a) gives $2$. The claim also
refutes the literal formula of cote 44, scan 4, under just its recorded
geometrically integral-fibre hypothesis.

\textbf{Proof.}
Write
\[
 X=\{uv+\textstyle\sum_{i=1}^{m} w_i z_i=0\}
       \subset\mathbf P^{2m+1}.
\]
Use the pencil of hyperplanes $a u+b v=0$, with parameter $[a:b]$.
Its base is the smooth quadric
\[
 B=\{u=v=0,\ \textstyle\sum w_i z_i=0\}\simeq Q^{2m-2}.
\]
The members with $ab\ne0$ are smooth odd-dimensional quadrics. This remains
true in characteristic two: after eliminating $v$, the equation is
$-t u^2+\sum w_i z_i=0$; the only possible projective singular point
from its partial derivatives would have all $w_i,z_i$ zero, and then
the equation forces $u=0$ since $t\ne0$. The two
remaining members, $u=0$ and $v=0$, are cones over $B$, each with precisely
one ordinary quadratic singularity, at the $v$- or $u$-vertex respectively.
For $m\geq2$ the base $B$ is a smooth integral positive-dimensional quadric,
so its cones are integral. The axis is transverse, and the total space
obtained by blowing up $B$ is smooth. At each singular point the local
equation is $\sum w_i z_i=0$ in $2m$ variables; its polar form is
nondegenerate also in characteristic two. Thus this is a Lefschetz pencil in
the precise sense of SGA 7 II, XVII, 2.2. It lies in the generic orbit of
pencils under the orthogonal group: equivalently its pencil line meets the
dual quadric transversely in two points.

Put
\[
 L_+=\mathbf P\langle w_1,\ldots,w_m\rangle,\qquad
 L_-=\mathbf P\langle w_1,\ldots,w_{m-1},z_m\rangle
       \subset B.
\]
These are maximal $(m-1)$-planes in the two different families on $B$.
Let $P_{\infty,+},P_{\infty,-}$ be their cones with the $u$-vertex and
$P_{0,+},P_{0,-}$ their cones with the $v$-vertex. They are $m$-planes on
$X$. Denote their Chow classes by the same symbols and set
\[
 Z=P_{\infty,+}-P_{\infty,-}.
\]
The two planes in this difference belong to the two different families on
$X$. Replacing the vertex while retaining $L_+$ interchanges these families;
the same holds for $L_-$. This assertion can be checked without a
classification theorem. Quotient by
$\langle w_1,\ldots,w_{m-1}\rangle$ and use the rank-one matrix description
of the quadric surface $uv+w_mz_m=0$, with matrix
$\left(\begin{smallmatrix}u&w_m\\-z_m&v\end{smallmatrix}\right)$.
The two row lines are fibres of one ruling and the two column lines are
fibres of the other. Varying the row (respectively column) gives a
$\mathbf P^1$-family joining $P_{\infty,+}$ to $P_{0,-}$ (respectively
$P_{\infty,-}$ to $P_{0,+}$), after adjoining the common subspace.
Consequently
\[
 [P_{0,+}]-[P_{0,-}]=-Z.
\]

We check the intersection number directly. For an $m$-plane $P\subset X$,
the normal-bundle sequence is
\[
 0\longrightarrow N_{P/X}\longrightarrow
 \mathcal O_{\mathbf P^m}(1)^{\oplus(m+1)}
 \longrightarrow\mathcal O_{\mathbf P^m}(2)\longrightarrow0.
\]
The last map is surjective since $X$ is smooth. Therefore
\[
 [P]^2=[h^m]\frac{(1+h)^{m+1}}{1+2h}
       =\frac{1+(-1)^m}{2}.
\]
Here $[h^m]$ denotes extraction of the coefficient of $h^m$. If $\alpha$
and $\beta$ are the two middle plane classes, a linear section degenerating
to two planes gives $h^m=\alpha+\beta$. Since $\deg(h^m\alpha)=1$,
\[
 \alpha\beta=\frac{1-(-1)^m}{2},\qquad
 (\alpha-\beta)^2=2(-1)^m.
\]
For the degeneration $h^m=\alpha+\beta$, intersect $X$ with the
$\mathbf P^{m+1}$ on which all $z_i$ vanish. Its equation is $uv=0$, and
its components are $P_{\infty,+}$ and $P_{0,+}$; the row--column families
above identify $[P_{0,+}]=\beta$. Thus every Chow relation used in this
calculation has an explicit geometric realization.

On the generic fibre, take the affine parameterization $v=-t u$. Its
equation is $-t u^2+\sum w_i z_i=0$. The restrictions of the two planes
appearing in $Z$ are $L_+$ and $L_-$. They are rationally equivalent over
$K$, not merely after passing to $\overline K$. Indeed, put
$W=\langle w_1,\ldots,w_{m-1}\rangle$ and consider the smooth $K$-conic
\[
 C_t=\{-t u^2+w_mz_m=0\}\subset\mathbf P^2_K.
\]
It is parameterized by
\[
 [r:s]\longmapsto[u:w_m:z_m]=[rs:r^2:t s^2].
\]
The family of projective $(m-1)$-planes
$\mathbf P(W\oplus\ell)$, for the isotropic lines $\ell$ represented by
$C_t$, lies in $Y$. At $[r:s]=[1:0]$ and $[0:1]$ its members are $L_+$
and $L_-$. Since these are $K$-rational points of $C_t\simeq\mathbf P^1_K$,
the two fibre cycles are rationally equivalent. Thus $Z|_Y=0$ as asserted.

We now compute the correction at each singular fibre. Write
$q:\widetilde X=\operatorname{Bl}_{B}X\to X$,
$j:E\hookrightarrow\widetilde X$, and $p:E\to B$. Set
\[
 t_X=q^*Z,\quad \delta_B=[L_+]-[L_-],\quad
 e=j_*p^*\delta_B,
\]
and let $u_\infty,u_0$ be the differences of the strict transforms of
$P_{\infty,\pm}$ and $P_{0,\pm}$, respectively. Then
\[
 u_\infty=t_X-e,\qquad u_0=-t_X-e.                 \tag{H9.1}
\]
For completeness, these identities follow directly from the incidence
model $\widetilde X\subset X\times\mathbf P^1$. On
$P_{\infty,\pm}\times\mathbf P^1$, the incidence equation is $a u=0$;
its divisor is the union of the plane in the fibre $a=0$ and
$L_\pm\times\mathbf P^1$. This is precisely the refined pullback of
$P_{\infty,\pm}$. On $P_{0,\pm}\times\mathbf P^1$ the equation is
$b v=0$, yielding the analogous identity. Taking differences proves
(H9.1), including the coefficients and signs.

The projection formula gives $t_X e=0$ because $q_*e=0$. Since
$\mathcal O_E(E)=\mathcal O_E(-1)$ and $p_*c_1(\mathcal O_E(1))=1$,
\[
 e^2=-\delta_B^2=-2(-1)^{m-1}=2(-1)^m,
 \qquad t_X^2=2(-1)^m.
\]
It follows that
\[
 \begin{array}{c|cc}
  &u_\infty&u_0\\\hline
 u_\infty&4(-1)^m&0\\
 u_0&0&4(-1)^m
 \end{array},\qquad
 t_Xu_\infty=2(-1)^m,\quad t_Xu_0=-2(-1)^m.
                                                        \tag{H9.2}
\]
In particular $t_X=(u_\infty-u_0)/2$ is itself a rational vertical cycle.
Its admissible correction is
\[
 t_X-\tfrac12u_\infty+\tfrac12u_0=0.
\]
Each local Schur-complement term is
$(2(-1)^m)^2/(4(-1)^m)=(-1)^m$. Thus each singular fibre contributes $1$
after multiplication by the H9 sign $(-1)^m$, and the two contributions
cancel the uncorrected value $2$. This proves the counterexample and the
claimed local refinement.\hfill$\square$

This is a mathematical insufficiency of the stated hypothesis, not evidence
that the word ``int\`egres'' was transcribed incorrectly. Without images of
the original manuscript no palaeographic alteration is justified.

\subsection*{H9-T1. The corrected intersection formula}

Let $f:\mathcal X\to C$ be a projective morphism from a smooth projective
$2m$-fold to a smooth projective curve over an algebraically closed field,
with smooth generic fibre $Y$. Use rational Chow groups. Call
$x\in\operatorname{CH}^m(\mathcal X)_{\mathbf Q}$ admissible if
$x|_Y$ is numerically trivial and
\[
 \deg(xv)=0
 \quad\text{for every codimension-$m$ cycle $v$ supported in fibres.}
                                                        \tag{H9.3}
\]
The geometric intersection-height normalization used here is
\[
 h_m(x|_Y,y|_Y)=(-1)^m\deg(xy)                    \tag{H9.4}
\]
for admissible extensions. Its independence of extensions follows from
the localization exact sequence: two extensions of the same Chow class
differ, modulo rational equivalence, by a vertical cycle, and (H9.3)
annihilates this difference. This is the signed degree of the pairing in
Kahn, \emph{Refined height pairing}, Definition 2.2, Proposition 2.8 and
equation (2.9). The terminology ``geometric intersection height'' specifies
the construction; comparison with a particular Bloch biextension or a
N\'eron--Tate pairing is a separate compatibility statement.

Suppose $x_0,y_0$ extend numerically trivial generic-fibre cycles and that
vertical corrections are to be made inside a finite-dimensional space
$W$ spanned by numerical classes of vertical cycles. Let
$R=\operatorname{rad}(W)$ for the intersection form. Assume $x_0$ and $y_0$
are orthogonal to $R$. Choose representatives $v_1,\ldots,v_r$ for a basis
of $W/R$, put $M_{ij}=\deg(v_iv_j)$ and
\[
 a_i=\deg(x_0v_i),\qquad b_i=\deg(y_0v_i).
\]
Then $M$ is invertible, the corrected classes are
\[
 x=x_0-\sum_i(M^{-1}a)_i v_i,\qquad
 y=y_0-\sum_i(M^{-1}b)_i v_i,
\]
and direct bilinear expansion gives their signed intersection. If the
corrected classes are admissible (in particular, if $W$ contains every
vertical numerical class), this is the height:
\[
 h_m(x|_Y,y|_Y)=(-1)^m\left(\deg(x_0y_0)-a^{\!t}M^{-1}b\right).
                                                        \tag{H9.5}
\]
The expression is independent of the chosen representatives modulo $R$.
For a smaller $W$, the same expression still computes
$(-1)^m\deg(xy)$, but is not called a height until orthogonality to the
remaining vertical classes has been verified.
If $W$ contains all the vertical numerical classes, the corrected cycles
satisfy (H9.3). Conversely, existence of a correction in $W$ requires
$x_0\perp R$, because adding an element of $W$ cannot change its pairing
with $R$. The finite-dimensional linear algebra proves this criterion is
also sufficient. Hence (H9.5) states both the formula and its exact
existence gate.

For a pencil with blowup map $q$ and $x_0=q^*Z$, $y_0=q^*Z'$, the
projection formula gives $\deg(x_0y_0)=\deg(ZZ')$. Fibres at different
closed points are disjoint, so the matrix can be organized into local
blocks, after separating the radical. Formula (H9.5) is the requested
local variant, including higher-dimensional vertical cycles; it is not
restricted to irreducible components of fibres. For $m=1$ it recovers the
usual intersection-matrix correction in Shioda, \emph{On the Mordell--Weil
lattices}, Theorem 8.6, equations (8.11)--(8.15), printed p.~228. The
original printed page was rendered and read in this run.

\subsection*{H9-T2. A cohomological hypothesis that removes all corrections}

\textbf{Claim.} Let a Lefschetz pencil on a smooth projective $2m$-fold
$X$ have generic fibre $Y$ and smooth blowup $q:\widetilde X\to X$.
Fix $\ell\ne\operatorname{char}k$. Suppose at each singular fibre $Y_s$
the specialization map
\[
 H^{2m}(Y_s,\mathbf Q_\ell(m))\longrightarrow
 H^{2m}(Y_{\bar\eta_s},\mathbf Q_\ell(m))       \tag{H9.6}
\]
is injective. If $Z,Z'$ restrict to homologically trivial cycles on $Y$,
then $q^*Z,q^*Z'$ are admissible and
\[
 h_m(Z|_Y,Z'|_Y)=(-1)^m\deg(ZZ').               \tag{H9.7}
\]
This statement holds in every characteristic subject to the explicit
injectivity hypothesis (H9.6).

\textbf{Proof.} On the smooth locus of the pencil the fibre classes of
$q^*Z$ form a section of the locally constant sheaf
$R^{2m}f_*\mathbf Q_\ell(m)$. The section is zero at the geometric generic
point, hence zero on this locus. At a singular fibre, compatibility of the
class of a global cycle with specialization and (H9.6) imply
$i_s^*\operatorname{cl}(q^*Z)=0$ in the ordinary cohomology of $Y_s$,
where $i_s:Y_s\hookrightarrow\widetilde X$. For any $m$-dimensional cycle $V$
inside any fibre, the compatibility of the cycle class with cap products
and the trace gives
\[
 \deg(q^*Z\cdot V)=
 \operatorname{Tr}_{V}(i_s^*\operatorname{cl}(q^*Z)\cap[V])=0.
\]
Thus $q^*Z$ is admissible; similarly for $Z'$. Formula (H9.7) follows
from (H9.4) and the projection formula.\hfill$\square$

In the ordinary nodal Picard--Lefschetz setting, a sufficient condition
for (H9.6) is that every local vanishing cycle is nonzero in the middle
cohomology of the nearby smooth fibre. Indeed, the relevant exact sequence
is
\[
 H^{2m-1}(Y_{\bar\eta_s})\longrightarrow\mathbf Q_\ell(-m)
 \longrightarrow H^{2m}(Y_s)
 \longrightarrow H^{2m}(Y_{\bar\eta_s})\longrightarrow0.
\]
The first arrow is pairing with the vanishing cycle. Poincar\'e duality
makes it surjective when that cycle is nonzero, proving injectivity of the
next specialization map. This identifies the additional local condition
missing in H9; the ordinary quadric pencil has zero middle cohomology and
therefore zero vanishing cycles. The precise SGA framework is SGA 7 II,
XVIII, 5.3.2 (printed p.~287), 5.3.4--5.3.5 (p.~288): condition $(A)$ says
that the direct-image sheaves agree with the extensions of their
restrictions to the smooth locus. Proposition 6.2 (p.~319) identifies
$(A)$ with nontrivial inertia action on middle cohomology at every
singular fibre. These original printed pages were independently read in
the programme audit. Using condition $(A)$, rather than an unqualified
quadratic-monodromy convention, makes (H9.6) available in every
characteristic covered by that condition.

Assuming the attached 032 Corollary 8.4, if $X$ is an abelian $2m$-fold,
the right side of (H9.7) is positive definite on its primitive algebraic
cohomology. If $h$ is the ample hyperplane class defining the pencil, put
\[
 C_{\rm prim}=\{Z\in\operatorname{CH}^m(X)_{\mathbf Q}:
                          h\operatorname{cl}(Z)=0\},\quad
 R=\operatorname{im}(C_{\rm prim}\to\operatorname{CH}^m(Y)_{\mathbf Q}),
 \quad
 R_0=\operatorname{im}(\operatorname{CH}^m_{\rm hom}(X)_{\mathbf Q}
                         \to\operatorname{CH}^m(Y)_{\mathbf Q}).
\]
Here $\ker(j^*:H^{2m}(X)\to H^{2m}(Y))=\ker L_h$: indeed
$j_*j^*=L_h$, and $j_*:H^{2m}(Y)\to H^{2m+2}(X)$ is injective by
Poincar\'e duality and weak Lefschetz in degree $2m-2$. Thus every element
of $C_{\rm prim}$ satisfies the restriction hypothesis of H9-T2.
For pencils satisfying (H9.6), $h_m$ on $R$ has radical exactly $R_0$:
the inclusion of $R_0$ in the radical follows because a homologically
trivial ambient class has zero intersection with every algebraic class;
the opposite inclusion follows from (H9.7) and positive definiteness of
the primitive form. In particular the restriction map induces an
isometry
\[
 C_{\rm prim}/\operatorname{CH}^m_{\rm hom}(X)_{\mathbf Q}
       \xrightarrow{\sim}R/R_0
\]
for the signed primitive intersection form on the left and the geometric
height on the right. For injectivity, if a restriction belongs to $R_0$,
subtract a homologically trivial ambient lift; the resulting restriction
is zero, so its height is zero and positivity forces its ambient
cohomology class to vanish. This proves the assertion without pretending
that a cohomology class has a canonical Chow representative. The argument
does not turn homologically trivial restrictions into points of an
algebraic intermediate Jacobian.

\subsection*{H9-T3. The precise quotient statement behind H9(b)}

Here is the valid linear-algebra implication, separating all geometric
inputs. Let $V$ be a rational vector space of cycle classes equipped with
a symmetric form $B$, and let
\[
 a:V\longrightarrow
 (J(K)/\operatorname{Tr}_{K/k}J(k))\otimes\mathbf Q
\]
be a defined homomorphism. Suppose $J$ is polarized and a height comparison
has established
\[
 B(v,w)=\langle a(v),a(w)\rangle_{\rm NT}.
                                                        \tag{H9.8}
\]
Then $B$ is positive semidefinite and
\[
 \operatorname{rad}B=\ker a.
\]
Indeed the right side is positive definite on the Lang--N\'eron space
after quotienting torsion, so $B(v,v)=0$ is equivalent to $a(v)=0$; the
bilinear identity proves the radical statement. Thus $B$ is positive
definite on $V$ if and only if $a$ is injective; it is always positive
definite on $V/\ker a$.

Applied to a pencil, (H9.8) must use the corrected form (H9.5). Replacing
it by the uncorrected primitive intersection form additionally requires
that every local correction vanish. The quadric example has nonzero
primitive classes with zero generic restriction, so injectivity on the
original primitive space fails. This does not by itself refute a bare
logical equivalence of two positivity assertions (both are true for a
quadric); it refutes the claimed identification and its asserted proof.

There is also a distinct domain obstruction. H9 assumes only homological
triviality, whereas Grothendieck's $P^m$ is constructed from algebraically
trivial cycles. SGA 7 II, XX, Theorem 4.1 and Remark 4.3 provide examples
where these restrictions are not algebraically trivial, even with
$X=E^4$: with $a_i$ the divisor class from the $i$th elliptic factor and
$h=\sum a_i$, the class
\[
 \alpha=(a_1-a_2)(a_3-a_4)
\]
satisfies $h\alpha=0$ and $\alpha^2=4$. On sufficiently high-degree
Lefschetz hyperplane sections, its restriction is homologically trivial
but not algebraically trivial. The original locators are printed
pp.~355--356; the full source hypotheses and their verification are
supplied by the independent programme audit. In particular, the regular
homomorphism from algebraically trivial cycles defining $P^m$ cannot be
applied to this cycle. One can ask for other, additional maps or
projections on a larger group, but H9 has not specified or proved such a
construction. Thus its displayed ``image'' in the algebraic $P^m$ of
cote 44 is not defined on the stated domain by the construction cited.

The repaired route must retain three pieces: the admissible height on
homologically trivial cycles; the local vertical form; and, only on the
algebraically representable part, comparison with an abelian
N\'eron--Tate height. The supplied CM-Hodge and standard-conjecture theorems
give positivity of the primitive intersection form for abelian $X$; they
do not remove either of the two explicit counterexamples above.

\subsection*{H9-T4. Polarization of the fixed primitive part of an ample section}

\textbf{Claim.} Assume 032, Corollary 8.4. Let $X$ be an abelian variety
of dimension $2m$ over an algebraically closed field of any characteristic,
let $j:Y\hookrightarrow X$ be a smooth ample divisor, and write
$h=[Y]$, $n=2m-1$. The fixed primitive summand of the incidence higher
Picard variety $P^m(Y)$ has a rational polarization given by its
$(-1)^m$-signed incidence correspondence. Its cohomological subspace is
\[
 F=j^*\ker\bigl(L_h^2:H^{n}(X)\longrightarrow H^{n+4}(X)\bigr)
       \subset H^n(Y).
\]
This proves the relevant part of Theorem 4 even though $Y$ need not be
abelian. The residual summand is the variable part, on which $j_*$ is zero.

\textbf{Proof.}
Hard Lefschetz on $X$ gives an isomorphism
$L_h:H^n(X)\to H^{n+2}(X)$; its inverse $\beta$ is algebraic because
the Lefschetz standard conjecture holds for abelian varieties. Let $e$
be the algebraic primitive projector in degree $n$ on $X$, with image
$\ker L_h^2$, and extend it by zero in other degrees. Define
\[
 r=\beta j_*,\qquad \Pi_F=j^*e\beta j_*.
\]
Since $j_*j^*=L_h$, one has $rj^*=1$ and $\Pi_F^2=\Pi_F$ on
cohomology. The correspondence $\Pi_F$ is algebraic. The incidence
detection lemma from JAC-T4.1 therefore induces an idempotent on
$P^m(Y)$ in the isogeny category.

The decomposition
\[
 H^n(Y)=j^*H^n(X)\oplus\ker j_*
\]
is orthogonal for the middle pairing on $Y$. On its first summand the
pairing is $\int_X uvh$. For $u\in\ker L_h^2$ and
$v=L_hw\in L_hH^{n-2}(X)$ one has
$\int_Xuvh=\int_X(L_h^2u)w=0$. Thus $\Pi_F$ is self-adjoint for
the pairing on $Y$. Saito's transpose compatibility and nondegenerate
incidence duality imply that its image $B_F$ is a nondegenerate
self-dual isogeny factor of $P^m(Y)$. Moreover it is primitive for $Y$:
if $u\in\ker L_h^2$, then
$j_*L_{j^*h}j^*u=L_h^2u=0$, and
$j_*:H^{n+2}(Y)\to H^{n+4}(X)$ is injective by dual weak Lefschetz.

It remains to verify the sign on all curves in $B_F$. Pull a normalized
rational Poincar\'e family back to a smooth pointed curve $C$ mapping to
$B_F$, and apply the parameter degree-one projector. Its class is an
algebraic tensor
\[
 \xi\in H^1(C)\otimes F.
\]
Set $\eta=(1\otimes r)\xi\in H^1(C)\otimes H^n(X)$.
It is algebraic, $j^*\eta=\xi$, and $h^2\eta=0$.
If $b$ is the genus of $C$ and $\theta$ the canonical principal
polarization on $J(C)$, the algebraic inverse to the Abel pullback
lifts $\eta$ to an algebraic class
$w\in H^1(J(C))\otimes H^n(X)$. It satisfies $h^2w=0$.
For $H=\theta+h$ on the abelian variety $J(C)\times X$,
\[
 H^{b+1}w=0.
\]
Indeed terms with $h$-exponent at least two vanish, while terms with
$h$-exponent zero or one have excessive degree on $J(C)$.
Consequently $w$ is primitive of codimension $m$. In its quadratic form
only one term of $H^b$ survives, and the Abel curve-class formula yields
\[
 \begin{aligned}
 \int_{J(C)\times X}w^2H^b
 &=b\int_{J(C)\times X}w^2\theta^{b-1}h\\
 &=b!\int_{C\times X}\eta^2h
 =b!\int_{C\times Y}\xi^2.
 \end{aligned}                                                \tag{H9.9}
\]
The factor $b!$ is positive. By 032, Corollary 8.4, the last intersection
has sign $(-1)^m$ and is zero only if $\xi=0$. The same curve test and
nef-plus-nondegenerate criterion as in JAC-T4.3 therefore prove that
the signed incidence class on $B_F$ is ample. This finishes the proof.
\hfill$\square$

The projectors used here implement the hyperplane formula of cote 16,
scan 40; the cohomological fixed/variable decomposition belongs to
SGA 7 II, XVIII, \S5. The input 032 is applied only to the genuinely
abelian product $J(C)\times X$, never to $J(C)\times Y$.

\subsection*{H9-T5. A finite-field Newton witness eliminates the remaining part}

\textbf{Claim.} Let a Lefschetz pencil of ample sections of an abelian
$2m$-fold be defined over a finite field of characteristic $p$. Write
$U$ for its smooth locus and $\mathcal V_\ell$ for its variable
$R^{2m-1}$-local system. Assume that $\mathcal V_\ell$ is nonzero and
remains absolutely irreducible under every finite-index restriction of
geometric monodromy, for a fixed prime $\ell\ne p$. Let
$B_{\rm var}$ be the variable incidence factor of $P^m(Y)$ over the
generic field, after a finite extension making the factor and universal
cycle defined. If one smooth closed fibre has a Frobenius eigenvalue on
its variable cohomology with normalized $p$-adic slope outside
$[m-1,m]$, then $B_{\rm var}=0$. In that case H9-T4 proves the primitive
Theorem 4 for the full incidence quotient of the generic fibre.

\textbf{Proof.}
The algebraic correspondence $1-j^*L_h^{-1}j_*$ induces the variable
idempotent on the incidence quotient. It is self-adjoint, so its image
$B_{\rm var}$ has nondegenerate incidence duality. A normalized Poincar\'e
cycle gives, for every $\ell\ne p$, an algebraic correspondence map
\[
 V_\ell B_{\rm var}\longrightarrow\mathcal V_\ell(m)
                                                        \tag{H9.10}
\]
over a finite cover of $U$. This map is injective: a vector in its
kernel pairs to zero with every other vector under the pullback of the
middle pairing, whereas that pullback is the realization of the
nondegenerate incidence isogeny. A choice of polarization on
$B_{\rm var}$ identifies $V_\ell B_{\rm var}$ with
$H^1(B_{\rm var},\mathbf Q_\ell)(1)$. Thus (H9.10) is equivalently
\[
 H^1(B_{\rm var},\mathbf Q_\ell)(-(m-1))
       \lhook\joinrel\longrightarrow\mathcal V_\ell.       \tag{H9.11}
\]
If $B_{\rm var}$ is nonzero, the irreducibility hypothesis, which survives
the finite field-of-definition extension, forces (H9.11) to be an
isomorphism.

At every point of the finite cover lying above a smooth fibre, the
right side is unramified. Therefore the left side is unramified and
the N\'eron--Ogg--Shafarevich criterion gives good reduction of
$B_{\rm var}$ there. Smooth proper base change specializes (H9.11).
If the residue field has cardinality $q$, the Frobenius eigenvalues on
$H^1$ of this reduction are algebraic integers $\gamma$ for which
$q/\gamma$ is also an algebraic integer. Hence, at every embedding into
$\overline{\mathbf Q}_p$,
\[
 0\leq\frac{v_p(\gamma)}{v_p(q)}\leq1.
\]
The Tate twist $-(m-1)$ multiplies geometric-Frobenius eigenvalues by
$q^{m-1}$. Thus all slopes on the variable part must belong to
$[m-1,m]$. Passing to a finite extension of a residue field raises
each eigenvalue and $q$ to the same positive power and leaves these
normalized slopes unchanged. A single smooth original fibre with a
slope outside the interval therefore contradicts $B_{\rm var}\ne0$.
This proves the claim.\hfill$\square$

The geometric irreducibility hypothesis is the standard monodromy
condition studied in SGA 7 II, XVIII, 6.7.2; it is stated explicitly
here so that no exceptional pencil is silently included. The good
reduction criterion and the integrality/duality of Frobenius on abelian
varieties are background inputs. This proof does not assert that every
high-degree pencil has the required Newton witness. Producing that
witness is the residual concrete gate of this method; purity of the
complex absolute values alone is insufficient to do so.
